\begin{lemma}[Situation 7.1: the maximal eligible index is penultimate]
Assume the Situation 7.1 data, including
\(\noderef{HierarchyPathNumber}=1\), and suppose the finite set
\(\noderef{HierarchyEligibleCommonMissingIndices}\) is nonempty.  Assume
also that every proper exactly-one-uncolored replacement coloring has a
missing color at every vertex.  Then
\[
 \ell-1\in\noderef{HierarchyEligibleCommonMissingIndices}.
\]
Thus some full replacement presentation has a color missing at both
\(y_{\ell-1}\) and \(y_\ell\).
\end{lemma}
\begin{proof}
Let \(I\) be the eligible-index finite set and choose its maximum \(i\);
this uses only finiteness of the interval filter, not a finite type of
replacement presentations.  Membership in the defining filter gives
\[
 1\leq i\leq\ell-1                                      \tag{3}
\]
and supplies an actual
\noderef{HierarchyReplacementPresentation} \(P\) and a color \(\alpha\)
missing at both \(y_i\) and \(y_\ell\).

Suppose \(i<\ell-1\).  Let \(D_{n,q}\) be the last
\noderef{HierarchyStoredColorsAt} of the good hierarchy carried by \(P\).
If \(\alpha\notin D_{n,q}\), apply
\noderef{HierarchySituationOneAdvanceOutsideStored} with (3), the same
presentation, and the same exact common-missing witnesses.  It gives
\(i+1\in I\).  If \(\alpha\in D_{n,q}\), apply
\noderef{HierarchySituationOneAdvanceInsideStored}; it gives the identical
conclusion.  In either case \(i+1>i\), contradicting maximality of \(i\).
Therefore \(i=\ell-1\), and rewriting the original membership proves the
claim.
\end{proof}
